\documentclass[a4paper,12pt]{article}
\usepackage[utf8]{inputenc}
\usepackage[english, russian]{babel}
\usepackage{amsmath, amssymb, amsthm}
\usepackage{enumitem}%оформление списков
\usepackage{varwidth}
\usepackage{graphicx}
\usepackage{float}
\usepackage{indentfirst}
\usepackage{url}
\usepackage{seqsplit}
\usepackage{pdflscape}
\usepackage{tabularx}
\usepackage{xcolor}
\usepackage{listings}
\usepackage{mathtools}
\usepackage{hyperref}
\usepackage[backend=biber, style=numeric, sorting=none]{biblatex}
\usepackage{titlesec}
\usepackage{appendix}
\usepackage{alltt}
\usepackage[hidelinks]{hyperref}
\usepackage{booktabs}

\addbibresource{my_libs.bib}


\usepackage[a4paper, top=20mm, left=30mm, right=20mm, bottom=20mm]{geometry}
\newcommand{\dim}{\mathrm{dim}}
\newcommand{\lin}{\mathrm{lin}}
\newcommand{\diag}{\mathrm{diag}}
\newcommand{\R}{\mathbb{R}}
\newcommand{\Rn}{\mathbb{R^n}}

\lstset{
    language=C++,
    basicstyle=\ttfamily\small,
    keywordstyle=\color{blue},
    commentstyle=\color{green!60!black},
    stringstyle=\color{red},
    numbers=left,
    numberstyle=\tiny\color{gray},
    stepnumber=1,
    breaklines=true,
    breakatwhitespace=true,
    tabsize=4,
    showspaces=false,
    showstringspaces=false
}

\newtheoremstyle{indented}% отступ
{3pt}% над
{3pt}% под
{\itshape}% шрифт тела, переключатель
{\parindent}% отступ
{\bfseries}% шрифт заголовка
{.}% точка после заголовка
{.5em}% отступ от заголовка
{}% head spec (empty = ‘normal’)

\theoremstyle{indented}
\newtheorem{theorem}{Теорема}
\newtheorem*{theorem*}{Теорема}
\newtheorem{definition}{Определение}
\newtheorem*{definition*}{Определение}
\newtheorem{following}{Следствие}
\newtheorem*{following*}{Следствие}
\newtheorem{statement}{Утверждение}
\newtheorem*{statement*}{Утверждение}

\renewcommand{\thesection}{\arabic{section}.}
\renewcommand{\thesubsection}{\arabic{section}.\arabic{subsection}.}
\renewcommand{\thesubsubsection}{\arabic{section}.\arabic{subsection}.\arabic{subsubsection}.}

\setcounter{section}{0}
\setcounter{subsection}{0}
\setcounter{equation}{0}

\usepackage{subfiles}

\begin{document}
\pagenumbering{arabic}
\setcounter{page}{1}
\pagestyle{empty}
\begin{center}
\textbf{Министерство науки и высшего образования Российской Федерации}
\\
\vspace{0.5cm}
\textbf{Федеральное государственное бюджетное образовательное\\
учреждение высшего образования} \\
\textbf{«Ярославский государственный университет им. П.Г. Демидова»}

\vspace{0.5cm}

Кафедра математического анализа
\vfill

\begin{flushright}
  Сдано на кафедру\\
  « \underline{\hspace{0.8cm}} » \underline{\hspace{3cm}} 2026 г.\\
  Заведующий кафедрой \\
  \underline{д. ф.-м. н., доцент }\\
  \underline{\hspace{3.4cm}} Невский М.В.\\
\end{flushright}
\vfill

Выпускная квалификационная работа магистра \\

\vspace{0.5cm}

\textbf{Задача о минимальной норме интерполяционного проектора} \\
(Направление подготовки магистров 01.03.02 Прикладная математика и информатика)
\vfill


\bigskip


\end{center}

\begin{flushright}
Научный руководитель\\
\underline{к. ф.-м. н., доцент }\\
\underline{\hspace{3.4cm}} Ухалов А.Ю.\\
« \underline{\hspace{0.8cm}} » \underline{\hspace{3cm}} 2026 г.\\
\vspace{1.5cm}
Студент группы \underline{ПМИ-21МО}\\
\underline{\hspace{3.5cm}} Зайков И.В. \\
« \underline{\hspace{0.8cm}} » \underline{\hspace{3cm}} 2026 г.\\
\end{flushright}

\vspace{0.5cm}
\begin{center}
Ярославль 2026 г.
\end{center}
\newpage

\section*{Реферат}
\addcontentsline{section}{Реферат}

Объем 52 страницы, 5 параграфов, 2 таблицы, 3 приложения. Список цитированной литературы содержит 6 источников.
\par
Изучается интерполяция непрерывной функции $f(x)$, определенной на кубе
$Q_n = [0, 1]^n$ , с помощью полиномов от $n$ переменных степени не выше 1. Ставится
задача оценивания минимальной нормы проектора для $n = 31$. Верхняя оценка $\theta_{31}$
получается из рассмотрения проектора с узлами интерполяции, совпадающими с
вершинами симплекса максимального объема в $Q_n$, построенного на основе матрицы Адамара. Получена новая оценка $\theta_{31} \leq 5.0$.
\par
\textbf{Ключевые слова: интерполяция, интерполяционный проектор, норма, матрица Адамара, симплекс максимального объема.}

\newpage

\pagestyle{plain}
\tableofcontents

\newpage
\subfile{main/0_main}
\subfile{main/1_code_structure}
\subfile{main/2_source_code_headers}

\end{document}